\subsection{11.3 Wirkung durch einen kontrollierten
Eingriff}\label{wirkung-durch-einen-kontrollierten-eingriff}

Ein nichtverschwindender Korrelator kann aus anfänglicher Überlappung
stammen. Ein aussagekräftiger Test vergleicht deshalb dieselbe Quelle
und denselben Eingang einmal mit und einmal ohne einen zugelassenen
Eingriff in A. Für gerade Observablen und \(U_A=e^{-i\varepsilon A}\)
gilt

\[\Delta_B(t)=\omega(U_A^\dagger B(t)U_A)-\omega(B(t))
=i\varepsilon\omega([A,B(t)])+O(\varepsilon^2).\]

Bei anfangs kommutierenden lokalen Bereichen ist der unmittelbare Effekt
null. Eine spätere Änderung kann einen kausalen Einfluss zeigen. Erst
eine hergeleitete Abstandsfunktion und kontrollierte Ausbreitung machen
daraus räumlichen Transport. Bei überlappenden Bereichen muss direkter
gemeinsamer Zugriff zuerst getrennt werden. Null lineare Antwort auf
einem einzelnen Zustand beweist umgekehrt nicht das Fehlen aller
Signale.

Ein jüngster exakter Gegenbeleg macht die Unterscheidung besonders klar.
Für zwei ungekoppelte Qubits mit

\[H=\tfrac12(Z_A+Z_B),\qquad
|\Omega\rangle=(|01\rangle+|10\rangle)/\sqrt2,\qquad A=X_A,\quad B=Y_B\]

ist der Zustand stationär und dennoch
\(\langle\Omega|B e^{-itH}A|\Omega\rangle=\sin t\). Gleichzeitig gilt
\([A,B(t)]=0\) für alle Zeiten. Die Kreuzkorrelation beginnt bei null
und verändert sich, während jedes unbedingte lokale spurtreue Instrument
in A die Statistik in B unverändert lässt. Dieses Gegenbeispiel wurde
hier symbolisch nachgerechnet. Eine Konditionierung auf einen A-Ausgang
ist eine andere Operation und kein Signal ohne Weitergabe dieses
Ausgangs. {[}S074, B3{]}

Lieb--Robinson-Abschätzungen liefern unter Lokalitäts- und Normannahmen
einen effektiven Ausbreitungskegel. Ihre Voraussetzungen müssen im
TFPT-Modell nachgewiesen werden; die Schranke allein erzeugt weder
Lorentzsymmetrie noch drei Raumdimensionen. {[}E6{]}

\subsection{11.4 Weshalb die Raumdimension offen
bleibt}\label{weshalb-die-raumdimension-offen-bleibt}

Aus demselben inneren Baustein kann man zunächst viele räumliche
Wiederholungen definieren. Ein gewählter Graph ist eine zusätzliche
Kompositionsregel. Weder seine Existenz noch die innere Zahl fünf wählt
automatisch ein dreidimensionales Kontinuum aus. Die Quellen enthalten
ausdrücklich alternative Verklebungen mit verschiedenen Dimensionen.

Ein zweikomponentiger hermitescher Operator besitzt drei
Pauli-Richtungen. Das erklärt die Form eines möglichen Weylsymbols
\(\sum_i v_iq_i\sigma_i\), sobald ein geeigneter Impulsraum vorhanden
ist. Es beweist nicht dessen Herkunft. Die bekannte Theorie von
Quantenautomaten gewinnt Weyl- und Dirac-Grenzformen unter konkret
festgelegten Voraussetzungen an Unitarität, Lokalität, Homogenität,
Isotropie und Graphstruktur. Diese Voraussetzungen dürfen nicht durch
die bloße Wiedererkennung der Pauli-Matrizen ersetzt werden. {[}E7{]}

\section{12. Relativistisches Feldwörterbuch, Chiralität und
Gravitation}\label{relativistisches-feldwuxf6rterbuch-chiralituxe4t-und-gravitation}

\subsection{12.1 Derselbe Matrixraum trägt eine
Lorentz-Kegelansicht}\label{derselbe-matrixraum-truxe4gt-eine-lorentz-kegelansicht}

Eine hermitesche Zweiermatrix

\[X=tI+x\sigma_x+y\sigma_y+z\sigma_z\]

hat Determinante

\[\det X=t^2-x^2-y^2-z^2.\]

Für \(A\in SL(2,\mathbb C)\) erhält \(X\mapsto AXA^\dagger\) diese
Determinante und den positiven Kegel. Das ist eine exakte
Lorentzstruktur auf diesem Darstellungsraum. Die Verbindung zum kleinen
Matrixkern ist deshalb mathematisch interessant. Sie identifiziert aber
noch keine dynamischen Raumzeitereignisse, keine gemeinsamen
Lichtgeschwindigkeiten verschiedener Sektoren und keine gravitative
Wirkung. {[}S046, Kap. 11, 34--35{]}

\subsection{12.2 Der Grassmann-Typentest bleibt
unverzichtbar}\label{der-grassmann-typentest-bleibt-unverzichtbar}

Sei \(M_A\) die antisymmetrische Fortsetzung einer W-Zeile und seien
\(\psi_{I\alpha}\) gleichhändige Weylfelder. Ein lokaler skalarer Ansatz
mit Gesamtkoeffizient \(M_A\otimes\epsilon_{\mathrm{Lorentz}}\) ist
symmetrisch: antisymmetrisch mal antisymmetrisch. Ein symmetrischer
Koeffizient kontrahiert mit zwei Grassmannfeldern zu null.

Ein symmetrischer Spinorkern liefert dagegen einen nichtverschwindenden
Typ \((1,0)\), ergänzt durch seinen konjugierten Typ. Das ist eine
Aussage über einen bestimmten gleichhändigen ableitungsfreien
Feldansatz. Der zusätzliche unabhängige Zweierfaktor
\(\epsilon_{\mathrm{aux}}\) kann einen skalaren Kanal retten, weil drei
antisymmetrische Faktoren insgesamt antisymmetrisch sind. Dieser Faktor
ist aber eine reale zusätzliche Ressource. Bei einer
Rang-eins-Projektion \(\psi_{Ia}=u_a\chi_I\) verschwindet er wieder:
\(\epsilon_{ab}u_au_b=0\). {[}S001, §5; S033, §6{]}

Die aktuelle innere Bilinearzerlegung ist vollständig:

\[\operatorname{End}(16\otimes4)
=(1,1)\oplus(45,1)\oplus(210,1)
\oplus(1,15)\oplus(45,15)\oplus(210,15).\]

Die Dimensionen \(1+45+210+15+675+3150=4096\) schließen den früher
unbestimmten Rest \(4035=210+675+3150\). Das zählt innere
Matrixrichtungen, keine 4035 neu hergeleiteten Teilchen. Die
Spinorkopplungen sind etablierte Darstellungstheorie. {[}S068, §5; E8{]}

\subsection{12.3 Algebraisch nicht null bedeutet noch nicht stabile
Feldtheorie}\label{algebraisch-nicht-null-bedeutet-noch-nicht-stabile-feldtheorie}

Der in S001 untersuchte freie geladene Zweitableitungsansatz für den
\((1,0)\)-Tensor besitzt eine nach unten unbeschränkte Energie. Dieser
konkrete Kandidat scheidet ohne zusätzliche Constraints aus. Daraus
folgt kein Verbot sämtlicher Tensorfeldtheorien.

Als Vergleich kann ein massiver Vektorparent mit Proca-Kinetik positive
freie Energie besitzen:

\[H_{\mathrm P}=\tfrac12\int d^3x\left[
\Pi^2+B^2+m^2A^2+\frac{(\nabla\cdot\Pi)^2}{m^2}\right]\ge0.\]

Eine Kopplung seiner selbstdualen Feldstärke an den symmetrischen
Paarstrom ist algebraisch möglich. Sie benötigt jedoch einen
zusätzlichen effektiven Vertrag und eine Skala. Freie Positivität
beweist weder wechselwirkende Positivität noch eine kontrollierte
Reduktion auf genau \(H_{\mathrm{nat}}\). {[}S001, §5{]}

Der notwendige Abschluss ist ein Adapter, der Ladung, Antisymmetrie,
Adjungierte, positive Kinetik und native Niedrigenergiedynamik zugleich
erhält. Unterschiedliche Feldtypen können parallel untersucht werden;
ihr Erfolg muss an derselben Quelle und denselben beobachtbaren
Prozessen gemessen werden.

\subsection{12.4 Familien und
Spiegelmoden}\label{familien-und-spiegelmoden}

Der innere Träger kann Familienzahlen und Anomaliesummen motivieren.
Drei physikalische Familien benötigen jedoch einen chiralen Operator
samt Darstellung, Maß und kontrollierter Spiegelentkopplung. Eine
Faltung oder Umbenennung von Weylknoten entfernt keinen zusätzlichen
dynamischen Freiheitsgrad. Im dokumentierten Walk treten mehrere
Weylknoten auf; ihre Chiralitäten und Quasienergien müssen im
tatsächlichen reziproken Gitter gezählt werden. {[}S046, Kap. 35{]}

Die Zuordnung \((16,4)\) ist außerdem keine automatische Gleichsetzung
des inneren Viererfaktors mit vier Lorentzkomponenten. Innere und
Raumzeitsymmetrie sind getrennte Transformationsgesetze. Gerade hier
verhindert eine sauber typisierte gemeinsame Konstruktion falsche
Identifikationen.

\subsection{12.5 Gravitation}\label{gravitation}

Eine Metrik auf Zustandsparametern, eine Lorentz-Kegelform oder ein
transversaler spurfreier Projektor ist noch kein dynamisches Graviton.
Erforderlich sind ein positiver masseloser Spin-2-Sektor, zwei physische
Helizitäten und eine konsistente universelle Kopplung im selben
dynamischen Modell.

In einer freien gapped Theorie beginnt die bilineare Spektralantwort
oberhalb einer Zweiteilchenschwelle; ein Projektor verschiebt diese
Energie nicht zu null. Eine wechselwirkende kritische Phase könnte
anders sein, müsste aber ausdrücklich konstruiert werden.
Entanglement-thermodynamische Einstein-Ableitungen setzen ihrerseits
wesentliche Eigenschaften einer lokalen QFT und eines Vakuums voraus.
Sie können diese Voraussetzungen nicht gleichzeitig als eigenes Ergebnis
ausgeben. {[}S046, Kap. 27; S071{]}

Die Projektorgeometrie aus Abschnitt 10 bietet deshalb eine mögliche
gemeinsame Sprache für Geometrie und innere Antwort. Sie ist kein Ersatz
für die dynamische Spin-2-Aufgabe.

\section{13. Physikalische Zahlen: Vorhersage, Kalibrierung und
gemeinsamer
Test}\label{physikalische-zahlen-vorhersage-kalibrierung-und-gemeinsamer-test}

Die frühen TFPT-Arbeiten enthalten zahlreiche Beziehungen für
Kopplungen, Massenskalen und kosmologische Größen. Diese gehören in ein
vollständiges Forschungspaper, müssen aber nach ihrem Transfervertrag
gelesen werden. Ein präziser dimensionsloser Formelwert ist erst nach
Zuordnung zur physikalischen Messgröße und Kontrolle der theoretischen
Unsicherheit eine physikalische Vorhersage.

\subsection{13.1 Feinstrukturkonstante}\label{feinstrukturkonstante}

Die aktuelle Compiler-Synthese berichtet

\[\alpha^{-1}_{\mathrm{Formel}}
=137.0359992168407125\ldots\]

Die am Recherchetag abgerufene NIST-Tabelle enthält weiterhin die
ausdrücklich bezeichnete CODATA-2022-Anpassung

\[\alpha^{-1}_{\mathrm{CODATA\,2022}}=137.035999177(21).\]

Der reine Abstand beträgt ungefähr \(1.90\) der dort angegebenen
Standardunsicherheit. Das ist ein deskriptiver Vergleich zwischen
Formelwert und datierter Referenz. Ohne unabhängige
Theoriefehlerkontrolle und vollständigen Thomson-Transfer ist es weder
eine allgemeine Bestätigung noch eine Widerlegung der TFPT. Die
Formelberechnung selbst wurde hier nicht als neuer Durchbruch
reproduziert. {[}S072, N6; E9{]}

\subsection{13.2 Kosmologie und gemeinsame
Parameter}\label{kosmologie-und-gemeinsame-parameter}

Für die im Hauptdokument untersuchte einfache Inflationsbranche gilt
nach Eliminieren der E-Faltungszahl

\[A_s(1-n_s)^2=\frac{c_3^7}{6\pi^2},\qquad
r=3(1-n_s)^2.\]

Damit sind Amplitude, Spektralindex und Tensorverhältnis innerhalb
dieser Näherung gekoppelt. Wer \(N\) aus \(A_s\) kalibriert, darf
\(A_s\) nicht anschließend als unabhängige Vorhersage zählen. Ebenso
dürfen dieselben Parameter nicht für verschiedene Observablen still
getrennt nachjustiert werden.

S046 dokumentiert einen datierten Branchentest gegen eine bestimmte
Kombination von ACT-DR6-Daten und eine Spannung in \(n_s\). Die konkrete
gemeinsame kosmologische Likelihood wurde für dieses Paper nicht neu
gerechnet. Der belastbare Schluss bleibt die Notwendigkeit eines
gemeinsamen Tests mit Reheating, Korrekturen, Datensatzwahl und
Theorieunsicherheit. Eine vereinfachte Formelbranche ist nicht die
gesamte TFPT. {[}S046, Kap. 27, 32--33{]}

\subsection{13.3 Was ein überzeugender empirischer Anschluss enthalten
muss}\label{was-ein-uxfcberzeugender-empirischer-anschluss-enthalten-muss}

Ein prüfbarer Ausgabevertrag fixiert vor dem Vergleich: denselben
Hamiltonoperator beziehungsweise dieselbe effektive Wirkung, dieselbe
Zustandsregel, eine endliche Menge unabhängiger Parameter, die Abbildung
auf Messgrößen, systematische Näherungsfehler und die verwendete
Datenversion. Werden Eingaben kalibriert, müssen die verbleibenden
unabhängigen Vorhersagen ausgewiesen werden.

Die native Kopplung \(g/\Delta\) ist bisher kein abgeleiteter Wert der
elektromagnetischen Kopplung. Der Polabstand \(\epsilon/\Delta\) ist
kein ausgewiesenes Elektronenmassenverhältnis. Gerade die Verbindung
dieser Größen mit den älteren Compilerformeln ist eine eigenständige
Aufgabe T6.

\section{14. Arithmetik, RH und
Rechenfähigkeit}\label{arithmetik-rh-und-rechenfuxe4higkeit}

\subsection{14.1 Der gemeinsame Begriff „Spektrum`` genügt
nicht}\label{der-gemeinsame-begriff-spektrum-genuxfcgt-nicht}

Ein Hamiltonspektrum, ein endlicher Markovoperator und eine
arithmetische explizite Formel sind verschiedene mathematische Objekte.
Eine Brücke zur Riemannschen Vermutung benötigt eine exakte, phasen- und
normierungstreue Identität mit dem vollständigen erforderlichen
Weil-Ausdruck oder ein anderes vollständig bewiesenes RH-Kriterium.
Positivität einiger Matrizen, einzelne Spektrallinien und die lokale
Focklücke reichen dafür nicht.

Ein allgemeiner positiver Prozesskern aus Abschnitt 8 wird nicht dadurch
zur RH-Lösung, dass man seine Indizes „Primschleifen`` nennt. Die
Gleichheit mit der arithmetischen Zielgröße muss einschließlich
archimedischer Beiträge, Kreuzterme, Randterme, Testraum und globalem
Grenzübergang bewiesen werden. Ein bereits positiv vorausgesetzter Kern,
dessen Identifikation mit der Weilform offen bleibt, verschiebt genau
die Hauptaufgabe.

\subsection{14.2 Grenze der aktuellen
Quellenprüfung}\label{grenze-der-aktuellen-quellenpruxfcfung}

Der aktuelle RH-Quellencheck stoppte mit \texttt{SOURCE\_UNAVAILABLE}
für den referenzierten Forschungsordner
\texttt{2026-09-04/scha-2/research}. Der bestehende Index enthält
hilfreiche routebezogene Hinweise, aber seine vollständige Aktualität
konnte dadurch nicht bestätigt werden. Die Unterbrechung ist kein
mathematischer Gegenbeweis; sie verhindert eine Behauptung, sämtliche
arithmetischen Arbeitsstände seien frisch vollständig geprüft.

Die aktuellen verfügbaren Quelltexte belassen die globale signierte
Weil-Positivität und kompatible kofinale Fortsetzung als offene
Bedingungen. In dieser Synthese wurde kein neuer RH-Kandidat als
bewiesen registriert und kein zentraler Beweismarker verändert. Der
eigenständige Forschungsschwerpunkt lag auf den physikalischen
Prozessverbindungen.

\subsection{14.3 Faktorisierung und
Universalität}\label{faktorisierung-und-universalituxe4t}

Ein kohärenter Gattervertrag kann algorithmische Universalität besitzen.
Das ist nicht identisch mit der physischen Verfügbarkeit seiner
Register, Kontrolloperationen, Resets und Messungen. Ein klassischer
Simulator eines Quantenschaltkreises hat auch dann keine automatisch
polynomiellen klassischen Kosten, wenn der ideale Quantenschaltkreis
eine effiziente Faktorisierung beschreibt.

Der Universalraum müsste für Rechenfähigkeit eine konkrete
Eingabe-Ausgabe-Ausführung samt Ressourcen, Fehlern,
Auslesewahrscheinlichkeiten und erfolglosen Versuchen liefern. P versus
NP wird weder durch Quantenfaktorisierung noch durch einen großen
Hilbertraum entschieden. Eine physikalische TOE müsste umgekehrt auch
nicht automatisch jedes offene mathematische Komplexitätsproblem lösen.

\section{15. Die fundamentale Gesamtlösung als präziser
Forschungsauftrag}\label{die-fundamentale-gesamtluxf6sung-als-pruxe4ziser-forschungsauftrag}

\subsection{15.1 Was sich auf einen gemeinsamen Gegenstand reduzieren
lässt}\label{was-sich-auf-einen-gemeinsamen-gegenstand-reduzieren-luxe4sst}

Die eigenen Herleitungen legen folgende Organisation nahe:

\[
\boxed{\text{Quellregel}
\longrightarrow (\mathfrak A,\omega,\Gamma,\alpha_t,\mathfrak I)
\longrightarrow\text{lokale und spektrale Geometrie}
\longrightarrow\text{physikalischer Grenzwert}.}
\]

Die erste Ebene legt zulässige Handlungen und ihre Zusammensetzung fest.
Der positive Wortkern konstruiert daraus den gemeinsamen erreichbaren
Raum. Seine Symmetriezerlegung trennt innere Typen von ihren Vorkommen.
Referenzsysteme bilanzieren erlaubte Ladungsänderungen. Dynamische
Kompression kontrolliert, welche niedrigen Antworten dieselbe
Entwicklung approximieren. Bewegte physische Projektoren können
geometrische Phasen und Metriken gemeinsam tragen.

Diese Kette verbindet viele bisher getrennte Themen mit wenigen
wiederkehrenden mathematischen Operationen: positiver Gramkern,
Quotientierung von Nullrichtungen, Isometrie, Komposition und spektrale
Projektion. Sie ist als mathematischer Zusammenhang konstruiert. Ihre
erste physische Auswahl und ihr letzter Grenzübergang sind noch nicht
gegeben.

\subsection{15.2 Warum „maximale Einfachheit`` allein nicht
auswählt}\label{warum-maximale-einfachheit-allein-nicht-auswuxe4hlt}

Ein Minimalitätsprinzip benötigt eine Vergleichsklasse und ein Maß für
Einfachheit. Ein minimaler Hilbertraum kann verschiedene Kernelphasen
tragen. Ein festes Symmetriealphabet lässt verschiedene
Kopplungsverhältnisse zu. Ein minimaler Referenzträger kann eine
mathematische Ladungsbilanz lösen und trotzdem keinen nativen Ursprung
haben. Ohne festgelegte physikalische Unterscheidbarkeit kann
Minimalität sogar beobachtbare Möglichkeiten versehentlich entfernen.

Die passende Forderung lautet deshalb: minimale Darstellung
\textbf{eines festgelegten vollständigen Prozessvertrags}, anschließend
Prüfung, ob die ursprüngliche Quelle diesen Vertrag eindeutig bis auf
operationelle Äquivalenz auswählt. Die GNS-Eindeutigkeit löst die erste
Frage unter gegebenem Kern. Die dimensionslose Momentenfamilie beweist,
dass die zweite Frage durch die bisherigen Symmetriedaten noch nicht
beantwortet ist.

\subsection{15.3 Drei priorisierte
Beweisaufgaben}\label{drei-priorisierte-beweisaufgaben}

\textbf{A. Ein ursprünglicher relationaler Übergang.} Gesucht ist ein
aus dem tatsächlich begründeten Alphabet erzeugtes Instrument, das den
neuen neutralen Kopplungstyp oder einen äquivalenten Gesamtprozess
realisiert. Zu liefern sind sämtliche inneren Ladungen, Referenzzustand,
Energie und Auslesung. Erfolg bedeutet Transfer auf dem gemeinsamen
Träger bei vollständiger Bilanz. Der oben bewiesene Ein-Cartan-Zeuge
dient als Zielvertrag, nicht als fertiger Ursprung.

\textbf{B. Ein physisch bewegter niedriger Projektor.} Auf demselben
Modell wird eine Familie \(P_x\) aus verfügbaren Operationen oder
dynamischen Freiheitsgraden hergeleitet. Zu prüfen sind \(dP_x\ne0\),
die Krümmung, eine Lücke oder kontrollierter Ersatz sowie die
tatsächlichen Messwahrscheinlichkeiten. Reine Rahmenrotationen und bloße
Nachselektion ohne Ausbeute bestehen diesen Test nicht. Die Parameter
müssen später aus operational bestimmten Teilen stammen, wenn eine
Raumzeitinterpretation beansprucht wird.

\textbf{C. Ein gemeinsamer Feld- und Größenadapter.} Derselbe geladene
Prozess muss in einem positiven, kausalen relativistischen Feldvertrag
und einer wachsenden Familie kontrolliert wiederkehren. Dabei sind
Ladung, Grassmannsymmetrie, Restantwort und sämtliche Übergangsfehler
gemeinsam zu erhalten. Erst diese Konstruktion trägt Aussagen über
Chiralität, Kopplungen, Kontinuum und Gravitation.

Als unmittelbar bearbeitbare Vorstufe dienen die gemischten
Momentenblöcke aus §6.4 und der Vergleich der Z4-Familie modulo der
nachgewiesenen Variablenwechsel aus §5.6. Ein Kodierungsnetz wäre
zusätzlich am Modenverlusttest aus §5.5 zu messen. Diese Aufgaben
liefern konkrete Entscheidungen über dieselbe Quelle, statt nur weitere
Spektralstellen zu sammeln.

Diese Reihenfolge ist eine Priorisierung, kein logisches Verbot
paralleler theoretischer Einsichten. Jede neue Lösung muss jedoch an
derselben Quelle anschließen. Genau dort trennt sich eine Verbindung von
einer bloßen Sammlung funktionierender Beispiele.

